\documentclass[10pt,a4paper]{amsart}
\usepackage[margin=19mm]{geometry}
\usepackage{amsmath,amssymb,mathtools}
\usepackage[scr=rsfs,cal=euler]{mathalfa}
\usepackage{xfrac}
\usepackage[unicode,hidelinks,bookmarks=false]{hyperref}
\newcommand{\ie}{i.e.\ }
\newcommand{\cf}{cf.\ }
\newcommand{\resp}{respectively\ }
\newcommand{\Subsection}{\subsection}
\providecommand{\nobreakdash}{\nobreak\hbox{-}}
\setlength{\emergencystretch}{2em}
\multlinegap0pt
\usepackage{mathtools}
\usepackage[shortlabels]{enumitem}
\usepackage{tikz}
\newtheorem{lemma}{Lemma}
\DeclareMathOperator*{\Argmin}{\MacroColor{argmin}}
\newcommand{\ConstraintSense}{\MacroColor{\ge}}
\newcommand{\EncodingSize}{size}
\newcommand{\FeasSolSet}{\MacroColor{\mathcal{F}}}
\newcommand{\MacroColor}[1]{#1}
\newcommand{\MyQED}{\mbox{}\hfill$\square$}
\newcommand{\NConstraints}{\MacroColor{m}}
\newcommand{\NVariables}{\MacroColor{n}}
\newcommand{\OneVector}{\mathbf{1}}
\newcommand{\OptSolSet}{\MacroColor{\mathcal{S}}}
\newcommand{\OptValFunc}{\MacroColor{v}}
\newcommand{\ReformPen}{\MacroColor{e}}
\newcommand{\Size}[1]{\langle#1\rangle}
\newcommand{\SizeThree}[3]{\langle#1,#2,#3\rangle}
\newcommand{\SizeTwo}[2]{\langle#1,#2\rangle}
\newcommand{\SolutionSizeEstimate}{\MacroColor{f_{\mathrm{sol}}}}
\newcommand{\StrictConstraintSense}{\MacroColor{>}}
\title{Complexity of Bilevel Linear Programming\\with a Single Upper-Level Variable}
\author{Nagisa Sugishita, Margarida Carvalho}
\date{Source: arXiv:2510.21126v2 (CC BY 4.0)}
\begin{document}
\maketitle

\noindent\emph{Source and licence.} Complexity of Bilevel Linear Programming\\with a Single Upper-Level Variable. Version arXiv:2510.21126v2 (CC BY 4.0), distributed by arXiv under the Creative Commons Attribution 4.0 International licence, http://creativecommons.org/licenses/by/4.0/, as recorded in the arXiv licence metadata for this exact version and shown on the abstract page https://arxiv.org/abs/2510.21126v2. Full licence notice: \texttt{licenses/arxiv-2510.21126-CC-BY-4.0.txt}.\par\smallskip

\noindent\textit{Editorial scope.} Counted unit: the article's lemma lemma:ZeroViolation (source0.tex lines 1240--1243). The article's complete original proof of it is delivered with it (source0.tex lines 1245--1353, one \texttt{proof} environment of 109 lines). No mathematics is written, repaired or re-derived: the statement and the proof are the article's own lines, in the article's own notation, and no retained line is substituted or re-flowed.  Retained uncounted as the source-local context the proof needs: lines 353--363; lines 370--387; lines 403--415; lines 418--439; lines 443--452.  Nothing was omitted from any retained span, so the package carries no omission record.  No retained line contains a citation, so the wrapper carries no bibliography and no bibliography entry.  No retained line contains a figure environment, an image-inclusion command or drawing code, so no image is delivered and the package declares no asset dependency.  Editorial formatting only, in addition: an AMS \texttt{amsart} preamble that restates the article's own declarations and macros used by the retained lines, namely the environments \texttt{lemma} on separate counters, together with the standard packages those lines need.  The retained environments print in this document as Lemma 1 in order of appearance, whereas the article prints the same items with the numbering of its own full document.  The complete native source of the article as distributed in the arXiv e-print for this exact version is bundled unchanged as \texttt{sources/187700/source0.tex}.
\begin{lemma}
\label{lemma:estimate-of-primal-linear-inequality}
Let $A, B \in \mathbb{Q}^{\NConstraints \times \NVariables}$, $a, b \in \mathbb{Q}^{\NConstraints}$, and $c \in \mathbb{Q}^{\NVariables}$.  
Let $\sigma$ denote the maximum \EncodingSize{} of the entries in $A$, $B$, $a$, and $b$.  
Consider the LP
\[
\min_{x} \left\{ c^{\top} x : A x \ConstraintSense a, B x \StrictConstraintSense b \right\}.
\]
If an optimal solution exists, then there exists a rational optimal solution $\bar{x}$ of size bounded by a polynomial in $\NVariables$ and $\sigma$.  
That is, there exists a polynomial $\SolutionSizeEstimate$ such that $\Size{\bar{x}} \le \SolutionSizeEstimate(\NVariables, \sigma)$.
\end{lemma}

\begin{lemma}
\label{lemma:SolutionSizeEstimate}
Let
$A_{1 1} \in \mathbb{Q}^{\NConstraints_1 \times \NVariables_1}, A_{1 2} \in \mathbb{Q}^{\NConstraints_1 \times \NVariables_2}, A_{2 1} \in \mathbb{Q}^{\NConstraints_2 \times \NVariables_1}, A_{2 2} \in \mathbb{Q}^{\NConstraints_2 \times \NVariables_2}$, $b_1 \in \mathbb{Q}^{\NConstraints_1}, b_2 \in \mathbb{Q}^{\NConstraints_2}$, $c_{1 1}, c_{2 1} \in \mathbb{Q}^{\NVariables_1}$ and $c_{2 2} \in \mathbb{Q}^{\NVariables_2}$.
Let $\NVariables = \NVariables_1 + \NVariables_2$ and $\sigma$ denote the maximum \EncodingSize{} of the entries in $A_{1 1}, A_{1 2}, A_{2 1}, A_{2 2}$, $b_1$, and $b_2$.
Consider the bilevel LP
\begin{align}
\min_{x_1, x_2} \ &
c_{2 1}^{\top} x_1 + c_{2 2}^{\top} x_2
\label{eq:SolutionSizeEstimateBLP}
\\
\text{s.t.} \ &
A_{2 1} x_1 + A_{2 2} x_2 \ge b_2, \notag \\
&
x_1 \in \Argmin_{x_1'} \{ c_{1 1}^{\top} x_1' : A_{1 1} x_1' + A_{1 2} x_2 \ge b_1 \}. \notag
\end{align}
If an optimal solution exists, then there exists a rational optimal solution $(\bar{x}_1, \bar{x}_2)$ such that $\SizeTwo{\bar{x}_1}{\bar{x}_2} \le \SolutionSizeEstimate(\NVariables, \sigma)$, \phantomsection\label{Rev:DefOfFSol} where $\SolutionSizeEstimate$ is the polynomial defined in Lemma~\ref{lemma:estimate-of-primal-linear-inequality}.
\end{lemma}

\begin{align}
\min_{x_1, x_2} \ &
c_{2 1}^{\top} x_1 + c_{2 2}^{\top} x_2 
\tag{$Q_{\mathrm{blp}}$}
\label{eq:PenaltyBefore}
\\
\text{s.t.} \ &
A_{2 1} x_1 + A_{2 2} x_2 \ge b_2, \notag \\
&
\bar{A}_{2 1} x_1 + \bar{A}_{2 2} x_2 \ge \bar{b}_2, \notag \\
&
x_1 \in \Argmin_{x_1'} \{ c_{1 1}^{\top} x_1' : A_{1 1} x_1' + A_{1 2} x_2 \ge b_1 \}, \notag 
\end{align}

\begin{align}
\min_{x_1, x_2, \ReformPen} \ &
c_{2 1}^{\top} x_1 + c_{2 2}^{\top} x_2 + M \ReformPen
\tag{$Q_{\mathrm{blp}}'(M)$}
\label{eq:PenaltyAfter}
\\
\text{s.t.} \ &
A_{2 1} x_1 + A_{2 2} x_2 \ge b_2, \notag \\
&
(x_1, \ReformPen) \in \Argmin_{x_1', \ReformPen'} 
\left\{ 
\begin{array}{rl}
&
A_{1 1} x_1' + A_{1 2} x_2 \ge b_1, \\
c_{1 1}^{\top} x_1' : 
&
\bar{A}_{2 1} x_1' + \bar{A}_{2 2} x_2 + \ReformPen' \OneVector \ge \bar{b}_2, \\
&
\ReformPen' \in I_{\ReformPen}
\end{array}
\right\}, \notag 
\end{align}

\begin{lemma}
\label{lemma:PenaltyReformulation}
Let $I_{\ReformPen}$ is $[0, 1]$ or $[0, \infty)$, and suppose \eqref{eq:PenaltyBefore} and \eqref{eq:PenaltyAfter} have an optimal solution for any $M > 0$.
Let $\NVariables = \NVariables_1 + \NVariables_2$, $\sigma$ denote the maximum \EncodingSize{} of entries in $A_{1 1}$, $A_{1 2}$, $A_{2 1}$, $A_{2 2}$, $\bar{A}_{2 1}$, $\bar{A}_{2 2}$, $b_1$, $b_2$, $\bar{b}_2$, and $c_{\infty} = \max\{ \| c_{2 1} \|_{\infty}, \| c_{2 2} \|_{\infty}, 1 \}$.
For any $M \ge 3 \NVariables c_{\infty} 4^{\SolutionSizeEstimate(\NVariables + 1, \sigma)}$, the following relations hold:
\begin{enumerate}
\item $\OptValFunc\eqref{eq:PenaltyBefore} = \OptValFunc\eqref{eq:PenaltyAfter}$;
\item $\OptSolSet\eqref{eq:PenaltyAfter} = \{(x_1, x_2, 0) : (x_1, x_2) \in \OptSolSet\eqref{eq:PenaltyBefore} \}$.
\end{enumerate}
\end{lemma}

\begin{lemma}
\label{lemma:ZeroViolation}
Under the assumptions of Lemma~\ref{lemma:PenaltyReformulation}, for any $(x_1, x_2, e) \in \OptSolSet\eqref{eq:PenaltyAfter}$, it holds that $e = 0$.
\end{lemma}

\begin{proof}
For any $x_1 \in \mathbb{R}^{\NVariables_1}$ and $x_2 \in \mathbb{R}^{\NVariables_2}$, $(x_1, x_2) \in \FeasSolSet\eqref{eq:PenaltyBefore}$ implies $(x_1, x_2, 0) \in \FeasSolSet\eqref{eq:PenaltyAfter}$.
Thus, we have 
\begin{equation}
\OptValFunc\eqref{eq:PenaltyBefore} \ge \OptValFunc\eqref{eq:PenaltyAfter}.
\label{eq:ProofPenaltyReformulationRelaxation}
\end{equation}
For the sake of contradiction, assume that there exists $(\bar{x}_1, \bar{x}_2, \bar{\ReformPen}) \in \OptSolSet\eqref{eq:PenaltyAfter}$ such that $\bar{\ReformPen} > 0$.
Under this assumption, we show that $\OptValFunc\eqref{eq:PenaltyBefore} < \OptValFunc\eqref{eq:PenaltyAfter}$, which contradicts \eqref{eq:ProofPenaltyReformulationRelaxation}.
Let $\NConstraints'$ be the number of constraints in the lower-level problem of~\eqref{eq:PenaltyAfter}.
Following the argument in the proof of Lemma~\ref{lemma:SolutionSizeEstimate}, we can write the feasible set of \eqref{eq:PenaltyAfter} as the union of polyhedra:
$$
\FeasSolSet\eqref{eq:PenaltyAfter} = \bigcup_{\omega \in \Omega'} U_{\omega}',
$$
where $\Omega' \subset \{0, 1\}^{\NConstraints'}$ and for each $\omega \in \Omega'$, $U_{\omega}' \subset \mathbb{R}^{\NVariables + 1}$ is a rational polyhedron defined by inequalities whose coefficients are of \EncodingSize{} at most $\sigma$.
Furthermore, there is $\bar{\omega} \in \Omega'$ such that $(\bar{x}_1, \bar{x}_2, \bar{\ReformPen}) \in U_{\bar{\omega}}'$.
Now, consider the LP
\begin{equation}
\min_{x_1, x_2, \ReformPen} \{ c_{2 1}^{\top} x_1 + c_{2 2}^{\top} x_2 + M \ReformPen : (x_1, x_2, \ReformPen) \in U_{\bar{\omega}}', \ReformPen > 0 \}.
\label{eq:PenaltyAfterStrictlyPositiveLP}
\end{equation}
Since $(\bar{x}_1, \bar{x}_2, \bar{\ReformPen})$ is also an optimal solution to $\eqref{eq:PenaltyAfterStrictlyPositiveLP}$, we have $\emptyset \not= \OptSolSet\eqref{eq:PenaltyAfterStrictlyPositiveLP} \subset U_{\bar{\omega}}' \subset \OptSolSet\eqref{eq:PenaltyAfter}$.
By Lemma~\ref{lemma:estimate-of-primal-linear-inequality}, there exists a rational solution $(\tilde{x}_1, \tilde{x}_2, \tilde{\ReformPen}) \in \OptSolSet\eqref{eq:PenaltyAfterStrictlyPositiveLP}$ such that $\SizeThree{\tilde{x}_1}{\tilde{x}_2}{\tilde{\ReformPen}} \le \SolutionSizeEstimate(\NVariables + 1, \sigma)$.
In particular, we have $\tilde{\ReformPen} > 0$, implying $\tilde{\ReformPen} \ge 1 / 2^{\SolutionSizeEstimate(\NVariables + 1, \sigma)}$.
Let $(\hat{x}_1, \hat{x}_2) \in \mathbb{Q}^{\NVariables}$ be a rational optimal solution to \eqref{eq:PenaltyBefore} such that $\SizeTwo{\hat{x}_1}{\hat{x}_2} \le \SolutionSizeEstimate(\NVariables + 1, \sigma)$.
Then,
\begin{align}
&
\OptValFunc\eqref{eq:PenaltyAfter}
-
\OptValFunc\eqref{eq:PenaltyBefore}
\notag
% \label{eq:LemmaPenatyReformulationContradiction}
\\
&
\qquad
=
(c_{2 1}^{\top} \tilde{x}_1 + c_{2 2}^{\top} \tilde{x}_2 + M \tilde{\ReformPen})
-(c_{2 1}^{\top} \hat{x}_1 + c_{2 2}^{\top} \hat{x}_2)
\notag
\\
% -------------------------------------------------
&
\qquad
\ge
M \tilde{\ReformPen} 
- \|c_{2 1} \|_\infty \| \tilde{x}_1 \|_1
- \|c_{2 2} \|_\infty \| \tilde{x}_2 \|_1
- \|c_{2 1} \|_\infty \| \hat{x}_1 \|_1
- \|c_{2 2} \|_\infty \| \hat{x}_2 \|_1
\notag
\\
% -------------------------------------------------
&
\qquad
\ge
M \tilde{\ReformPen} 
- c_\infty \NVariables_1 \| \tilde{x}_1 \|_\infty
- c_\infty \NVariables_2 \| \tilde{x}_2 \|_\infty
- c_\infty \NVariables_1 \| \hat{x}_1 \|_\infty
- c_\infty \NVariables_2 \| \hat{x}_2 \|_\infty
\notag
\\
% -------------------------------------------------
&
\qquad
\ge
M/2^{\SolutionSizeEstimate(\NVariables + 1, \sigma)}
- c_{\infty} \NVariables_1 2^{\SolutionSizeEstimate(\NVariables + 1, \sigma)}
- c_{\infty} \NVariables_2 2^{\SolutionSizeEstimate(\NVariables + 1, \sigma)}
\notag
\\
& 
\qquad
\qquad
\qquad
\qquad
\qquad
\qquad
-  c_{\infty} \NVariables_1 2^{\SolutionSizeEstimate(\NVariables + 1, \sigma)}
- c_{\infty} \NVariables_2 2^{\SolutionSizeEstimate(\NVariables + 1, \sigma)}
\notag
\\
% -------------------------------------------------
&
\qquad
\ge
M/2^{\SolutionSizeEstimate(\NVariables + 1, \sigma)}
- 2 c_{\infty} \NVariables 2^{\SolutionSizeEstimate(\NVariables + 1, \sigma)}
\notag
\\
% -------------------------------------------------
&
\qquad
\ge
M/2^{\SolutionSizeEstimate(\NVariables + 1, \sigma)}
- (2/3) M/2^{\SolutionSizeEstimate(\NVariables + 1, \sigma)}
\notag
\\
% -------------------------------------------------
&
\qquad
> 0,
\notag
\end{align}
which contradicts~\eqref{eq:ProofPenaltyReformulationRelaxation}.
Therefore, for any $(\bar{x}_1, \bar{x}_2, \bar{\ReformPen}) \in \OptSolSet\eqref{eq:PenaltyAfter}$ we have $\bar{\ReformPen} = 0$.
\MyQED
\end{proof}

\end{document}
